\scanpage{24}
\mspage{23}

$\widetilde N$, c'est un $\underline{O}_X$-$G$-module \deleted{\unsure{qui est contenu}} $\subset\inserted{\underline{\Gamma}^{\bullet}(F')=}\unsure{F'}$, et \deleted{contient} $\widetilde{\underline{\Gamma}^{\bullet}(F)}=F$. C'est donc le \deleted{\unsure{$F'$} cherché} faisceau cherché. Par construction, sa formation commute à \unsure{tt} changement plat $S'\to S$ de la base. Enfin, si $F$ est de type fini, \deleted{\unsure{alors $\exists$}} prouvons que son enveloppe $F'$ l'est. On peut supposer $S$ affine. Alors $\exists n$ tel que $F(n)$ \deleted{\unsure{pour $n$ assez gd}} engendré par ses sections, \deleted{soient alors}\inserted{donc par des} sections $\varphi_i$ \deleted{$\unsure{\in}$} $(1\leq i\leq N)$ $\in\Gamma(S,\underline{\Gamma}^n(F))$.

Considérons alors le sous-module \deleted{$\underline N^n$ de $\underline{\Gamma}^{\bullet}(E)$, engendré par les $\varphi_i$. Alors $\widetilde{\underline N}\subset F$ est en fait égal à $F$, comme on \unsure{const.} tt de suite. On peut donc répéter la construction \unsure{précédente}, en} $G$-invariant de $\underline{\Gamma}^{\unsure{n}}(F)$ engendré par les $\varphi_i$, soit $N^n$, et le sous Module de $\unsure{\underline E|E}$ qu'il définit. \uline{On vérifie} encore que c'est \uline{$F'$}, ce qui prouve que $F'$ est de type fini.....

\margin{Rédiger mieux \unsure{au début} de la démonstration, pour \unsure{pouvoir} \figure{\begin{tikzpicture}\draw[->] (0,0)--(.4,.45);\end{tikzpicture}} \unsure{y référer|la fois prochaine} \unsure{ici}.....}
